% corpus_supplement: Erdős problem #475
% source: https://www.erdosproblems.com/latex/475
% fetched_at: 2026-07-14T18:47:48Z
% payload_sha256: 6358e519ab70dca01cd7bdad313a2419c5fb9117969b36121c3c8dc9187cb067
% statement/remarks/references extracted by erdos_ingest.extract_source_page;
% rendered by erdos_ingest.render_tex — do not edit
\documentclass{article}
\usepackage{amsmath, amssymb, amsthm}
\usepackage{hyperref}
\usepackage{geometry}
\geometry{margin=1in}

\title{Erd\H{o}s Problem \#475}
\date{Last updated: 2026-02-23}
\author{Source: erdosproblems.com}

\begin{document}
\maketitle

\noindent\textbf{Status:} OPEN \quad \textbf{No prize}

\noindent\textbf{Tags:} number theory, additive combinatorics

\medskip
\noindent\textbf{Problem Statement:}

Let $p$ be a prime. Given any finite set $A\subseteq \mathbb{F}_p\backslash \{0\}$, is there always a rearrangement $A=\{a_1,\ldots,a_t\}$ such that all partial sums $\sum_{1\leq k\leq m}a_{k}$ are distinct, for all $1\leq m\leq t$?

\bigskip
\noindent\textbf{Remarks:}

A problem of Graham, who proved it when $t=p-1$. A similar conjecture was made for arbitrary abelian groups by Alspach. Such an ordering is often called a valid ordering.

This has been proved for $t\leq 12$ (see Costa and Pellegrini \cite{CoPe20} and the references therein) and for $p-3\leq t\leq p-1$ (see Hicks, Ollis, and Schmitt \cite{HOS19} and the references therein).

This has been proved for all sufficiently large primes; this is a consequence of four different types of result, each of which handles (using different methods) a different range of size of $A$.
{UL}
{LI} The small $A$ case: Kravitz \cite{Kr24} proved this for\[t \leq \frac{\log p}{\log\log p}.\](This was independently earlier observed by Will Sawin in a MathOverflow post.)
Bedert and Kravitz \cite{BeKr24} improved this to the range\[t \leq e^{c(\log p)^{1/4}}\]for some constant $c>0$, and Costa and Della Fiore \cite{CoDe26} have further improved this to\[t \leq e^{c(\log p)^{1/3}}.\]{/LI}
{LI} The medium $A$ case: Pham and Sauermann \cite{PhSa26} proved this, for any $0<\alpha<1$, for all\[ 1 \ll_\alpha t\leq p^{1-\alpha}.\]{/LI}
{LI} The large $A$ case: Bedert, Buci\'{c}, Kravitz, Montgomery, and M\"{u}yesser \cite{BBKMM25} proved this for all\[p^{1-c}\leq t\leq (1-o(1))p\]for some small constant $c>0$.{/LI}
{LI}The very large $A$ case: M\"{u}yesser and Pokrovskiy \cite{MuPo25} proved this for all\[t\geq (1-o(1))p.\]{/LI}
{/UL}

\bigskip
\noindent\textbf{References:}

[BBKMM25] B. Bedert, M. Buci\'{c}, N. Kravitz, R. Montgomery, and A. M\"{u}yesser, On Graham's rearrangement conjecture over $\mathbb{F}_2^n$. arXiv:2508.18254 (2025).
 
 [BeKr24] B. Bedert and N. Kravitz, Graham's rearrangement conjecture beyond the rectification barrier. arXiv:2409.07403 (2024).
 
 [CoDe26] S. Costa and S. Della Fiore, New bounds for (weak) sequenceability in $\mathbb{Z}_k$. arXiv:2602.19989 (2026).
 
 [CoPe20] Costa, S. and Pellegrini, M. A., Some new results about a conjecture by Brian Alspach. Arch. Math. (Basel) (2020), 479-488.
 
 [HOS19] Hicks, Jacob and Ollis, M. A. and Schmitt, John R., Distinct partial sums in cyclic groups: polynomial method and
constructive approaches. J. Combin. Des. (2019), 369-385.
 
 [Kr24] N. Kravitz, Rearranging small sets for distinct partial sums. arXiv:2407.01835 (2024).
 
 [MuPo25] A. M\"{u}yesser and A. Pokrovskiy, A random Hall-Paige conjecture. arXiv:2204.09666 (2025).
 
 [PhSa26] H. T. Pham and L. Sauermann, On Graham's rearrangement conjecture. arXiv:2602.15797 (2026).

\bigskip
\noindent\small{Source: \url{https://www.erdosproblems.com/475}}
\end{document}
